\scanpage{61}
formée des $\{\zeta'_i\ (i\in I-K),\ \zeta_J\ (J\in\mathfrak P_2(K)\subset\mathfrak P_2(I))\}$.

La base \deleted{\unsure{associée à|associée}} à $B^{ij}$ est $B^{ji}$, \deleted{de} \inserted{celle ass.} à $B^K$ est $B'{}^K$.

\uline{Corollaire 1}\quad $\operatorname{Card}(\operatorname{Bases}(\Delta))=72$ \inserted{$=8.9=2^3 3^2$}, \deleted{$\operatorname{Card}\operatorname{Aut}(\Delta)=72\mathbin.$}
\[
\operatorname{Card}\underset{\substack{\Vert\\\operatorname{Aut}(\Delta)}}{W}=72.6!=2^7 3^{\unsure{4}}5
\]

\uline{Corollaire 2} $W$ est transitif sur $\Delta$ (ens des sommets), et \deleted{\unsure{sur}} sur l'ens des couples de sommets \inserted{distincts} \inserted{non} liés. Si $(\delta,\delta')$ est un couple de sommets distincts non liés, $\exists!$ base $I$ telle que $\delta\in I$ \deleted{\unsure{et|\&}}, $\delta'\in I'$ base associée (donc $\delta$ et $\delta'$ se correspondent par la bijection can $I\simeq I'$)

\uline{Corollaire 2} La cat. des complexes cubiques $\Delta$ munie d'un él $J\in\mathfrak P_2(\Delta)$ formé de $2$ sommets non liés, est équivalente à la cat. des couples formés d'un ens à cinq éléments et d'un autre à deux él. (\unsure{ce dernier étant|ce dernier est|ces derniers étant} $J$).

\deleted{\unsure{partie particulière|partie liée de cardinal}}

\deleted{Notion de $i$-base, \inserted{(de partie libre de card. $i$)} ($1\le i\le6$) :}

\deleted{\uline{Prop}. \unsure{Pour} \inserted{les} \unsure{bibases} \unsure{parties} libres \inserted{\unsure{ordonnées|ordonnée}} à $3$ él, \unsure{bibases} \unsure{deux} à deux \unsure{qui ne sont liés|non liés|opérées par $W$}. Pour les telles $J$, l'ens $K$ des éléments \inserted{de $\Delta-J$} de \unsure{$\Delta$|$\widetilde\Delta$} liés à tt él de $J$ est de cardinal $6$, et $\exists!$ bibase $\{I'\cup I''\}\simeq I\times\varepsilon$ telle que $J\subset I'\cup I''$, \unsure{stable} \illegible{2} \inserted{et dans} \unsure{le|$K$} \illegible{1} \unsure{les} $2$ él $K=\{K',K''\}$,}

% slanted
\margin{\deleted{et alors $J$ est \deleted{et $*K$} $\unsure{j'j|i'j}$ \unsure{au|ou} $j\in I_1=K/\unsure{i|2}$}}
\deleted{L'ens \inserted{$J'$} des élémts de $\Delta$ liés \unsure{et} \illegible{1} à aucun \unsure{point} de $J$ est de cardinal $3$ \unsure{et tt} él de $K$ est non lié à exactemt $1$ élément de \deleted{\unsure{$J$}} $J'$, d'où une application $K\longrightarrow\deleted{\unsure{\alpha}}J'$. Si la bibase est associée à l'ens à $2$ él. $\varepsilon$, on a donc $K\longrightarrow J'\times\varepsilon$ (ou $K\longrightarrow I'\cup I''\longrightarrow\varepsilon$) qui est un iso.}
% slanted
\margin{\deleted{\uline{NB} la relation entre $J$ et $J'$ est symétrique}}
% archive sheet number: 60
